\subsection{Different descriptions of the same assignment}

A formula is useful because it gives values for many inputs at once. A table emphasizes selected values, while a graph displays the pairs $(x,f(x))$ geometrically.

\begin{examplebox}[A function described by a table]
The table
\[
\begin{array}{c|ccc}
x&1&2&3\\ \hline
f(x)&4&4&7
\end{array}
\]
defines a function on the finite domain $\{1,2,3\}$. No algebraic formula is required.
\end{examplebox}

\begin{definitionbox}[Graph]
The graph of a real function $f:D\to\mathbb R$ is the set of points
\[
\{(x,f(x)):x\in D\}.
\]
\end{definitionbox}

A graph allows us to read approximate values, zeros, signs, intervals of increase or decrease, and visible discontinuities. These observations are geometric statements about the same assignment encoded by the formula or table.

\begin{interpretationbox}[The vertical-line test]
A drawn relation represents a function of $x$ only when every vertical line meets it at most once. Two intersections would assign two different outputs to the same input.
\end{interpretationbox}

\begin{corebox}[Read before calculating]
Before manipulating a formula, determine what its graph and domain already tell you. Later, limits and derivatives will refine this visual reading rather than replace it.
\end{corebox}
